\subsection{Mitigation: Evaluation and Further Strategies}\label{sec:mitigationDiscussion}
\begingroup
\let\subsection\subsubsection
\subsection*{Evaluation of $M1$}

\textcolor{red}{\textbf{Placeholder, to be replaced once the fits are complete.} This subsection reports the posterior of $\delta_O$ on both fitted sets with its credible interval, the probability $\Pr(\delta_O>0\mid D)$ against the first leg of the rule, and the leave-one-out comparison of the configuration level with and without the overlap term against the second leg, followed by the three-way verdict and the convergence diagnostics that gate it. It then qualifies the estimate with the three confounds this design cannot separate: the token count $N=1+\lfloor(L-P)/S\rfloor$ moves with the stride, so overlap and sequence length are not independently identified; lowering $S$ changes the mixture of stride- and patch-derived candidates along the same axis; and the truncated context tail of \autoref{sec:geometries} is non-zero on exactly the six geometries that identify the stride branch. None of the three could create a mitigation, but any could inflate one.}

\subsection*{Further strategies}

Because the candidate set follows from $(f_s,P,S)$ alone, two further strategies follow from the same arithmetic; neither is tested here.

The first is to move the combs rather than to thin one of them. Changing $S$ translates the stride comb and changing $P$ translates the patch comb, each by a known amount, so a deployment that knows the spectral content carried can choose a geometry whose combs fall between the components that matter instead of on them. It is a design-time control rather than a repair: it relocates the frequencies at which recovery is worst, deterministically and before any forecast is made, and protects only the content known in advance.

The second is redundancy across geometries: two models whose combs fall at different frequencies would be blind, if at all, at different frequencies, and an ensemble that selects per frequency, using the observability assessment of \autoref{sec:ontology} as the switch, would leave no frequency served only by a model known to be blind there. It addresses the effect rather than relocating it, at a cost: two forecasters must be calibrated against each other, the switch becomes a component that can be wrong, and the claim is untested here.

$M1$ thins one branch, geometry selection moves both, and only redundancy covers what a single tokeniser cannot represent.
\endgroup

\subsection{Security Risk Assessment}\label{sec:securityAssessment}
\begingroup
\let\subsection\subsubsection
The forecaster sits inside a control loop in which its output drives battery dispatch and curtailment on a photovoltaic micro-grid, so a forecasting failure can become a control failure, and the property tested in this report has to be read as a security property and not only as an accuracy one.

Phase-locking is deterministic, attacker-independent and present at deployment time, a consequence of the tokenisation geometry rather than of an adversary or of the training data. Structural aliasing is the phenomenon that may then appear in the learned representation, and phase-locking marks the sites where, following Pagani et al., a blind spot is most expected, without excluding others.

The candidate set of blind frequencies follows from the configuration $(f_s,P,S)$ alone. Where an exposure would sit, and how it moves when a configuration changes, are known before anything is measured. That is enough to place a monitor and to bound what a mitigation could reach; it is not enough to claim that a loss occurs.

\subsection{Risk Management Framework}
The NIST AI Risk Management Framework\autocite{tabassiArtificialIntelligenceRisk2023} is the reference adopted for that reading: it separates the trustworthiness characteristics at stake here, validity and reliability, security and resilience, from the governance functions that turn evidence into a control, and it asks that known limitations be mapped, measured and managed rather than merely disclosed.

Its four functions assign responsibility and oversight (GOVERN), fix the deployment context and the pathways to harm (MAP), track the likelihood, severity and uncertainty of each risk (MEASURE), and select and monitor responses (MANAGE); they are continuous and mutually informing, so \autoref{tab:riskRegister} records the functions each threat engages rather than a one-way sequence.

\textcolor{red}{\textbf{Placeholder, to be replaced once the fits are complete.} The representational cost at a candidate frequency is decided by $e^{\theta_{\mathrm{lock}}}$, the propagation of that cost to the forecast by the paired recovery ratio of Model~A, the dependence on phase by $\sigma_\phi$, and the predictability of the site locations by $\kappa_S$ and $\kappa_P$. Every threat below is therefore stated conditionally on the estimand that would establish it, under rules and thresholds fixed in \autoref{tab:bayesDecisions} before any posterior is inspected. Report here, for each of the four estimands named above, the posterior median, the credible interval and the probability against its pre-registered threshold; then state which of the pathways of \autoref{tab:riskRegister} the evidence establishes, which it leaves inconclusive and which it refutes, and revise the residual column accordingly. An estimand whose fit fails a convergence gate is reported as not reportable, and a branch below the identification bar as not identified, rather than as a negative finding.}

Each row of \autoref{tab:riskRegister} names a pathway, the framework functions it engages, the estimand that would establish it, the control available and the residual exposure, which stays conditional until that estimand has been read.

\begin{table*}[!t]
    \centering
    \small
    \setlength{\tabcolsep}{2.25pt}
    \begin{tabular*}{\textwidth}{@{\extracolsep{\fill}}>{\raggedright\arraybackslash}m{2.50cm}>{\raggedright\arraybackslash}m{3.00cm}>{\raggedright\arraybackslash}m{1.70cm}>{\raggedright\arraybackslash}m{3.85cm}>{\raggedright\arraybackslash}m{3.25cm}>{\centering\arraybackslash}m{1.45cm}@{}}
        \toprule
        \textbf{Risk} & \textbf{Description} & \textbf{AI RMF} & \textbf{Evidence in this report} & \textbf{Control available} & \textbf{Residual} \\
        \midrule
        Adversarial masking & A periodic component injected on a candidate frequency is under-represented and escapes forecast-residual detection & \shortstack[l]{MAP\\MEASURE} & Precondition is $e^{\theta_{\mathrm{lock}}}>1$ (Model~B) with the site locations of $\kappa_S$ (Model~D2); the pathway itself needs the paired recovery ratio of Model~A. Not yet evaluated & Spectral monitor on the raw stream, upstream of tokenisation, at the enumerated candidate set & high \\
        \cmidrule(lr){1-6}
        Systematic blind spot with no adversary & A legitimate plant harmonic sits on a candidate frequency and is persistently under-represented & \shortstack[l]{MAP\\MEASURE\\MANAGE} & Candidate set is closed-form in $(f_s,P,S)$ and needs no fit; the cost at those sites is decided by $e^{\theta_{\mathrm{lock}}}$. Not yet evaluated & Choose $(P,S)$ against the known spectral content of the plant; \valueof{PartiallyObservable} degradation & medium \\
        \cmidrule(lr){1-6}
        Configuration drift & A change of sampling interval, patch size or checkpoint silently relocates the blind set & \shortstack[l]{GOVERN\\MAP} & $H3a$: relocation is predicted in closed form by $f_s/S$, and whether the sites follow it is decided by $\kappa_S$. Not yet evaluated & Recompute the candidate set from the knowledge base at every configuration change & low \\
        \cmidrule(lr){1-6}
        False assurance & Overlap is increased and treated as a fix & \shortstack[l]{GOVERN\\MANAGE} & $M1$ is decided by $\delta_O$ under a two-leg rule and neither leg is evaluated. Independently of any fit, the patch branch is invariant to $S$ by construction & Do not credit overlap as a control; record the unestablished status in the model card & high \\
        \cmidrule(lr){1-6}
        Defence placed inside the blind spot & Phase randomisation, or anomaly detection on forecast residuals & MANAGE & $H2$ is decided by $\sigma_\phi$ (Model~C): a deficit independent of phase makes randomising it useless. Not yet evaluated & Move detection upstream of tokenisation; treat residual-based detection as insufficient alone & medium \\
        \bottomrule
    \end{tabular*}
    \caption{AI RMF names the function of the framework\autocite{tabassiArtificialIntelligenceRisk2023} under which each item is handled. Residual is the exposure that survives the control in the adjacent column, judged qualitatively: high where no control available to an operator closes the exposure, medium where a control exists but is partial or requires design-time choices, low where the exposure is closed by a procedure that follows from the arithmetic. Each entry names the estimand specified in \autoref{sec:six} that decides it; no posterior is asserted here, since the fits are not yet available. The two entries carried at high residual are those for which no control available to an operator closes the exposure: Adversarial masking, because a detector placed after tokenisation inherits the blind spot it is meant to find, and False assurance, because the pre-registered mitigation cannot reach the patch branch by the arithmetic of \autoref{sec:eight} whatever the fits return.}
    \label{tab:riskRegister}
\end{table*}


\subsection{Adversarial-masking threat model}
A weakness becomes a vulnerability when someone can reach it, and because phase-locking is deterministic, an adversary needs no model access, no query budget and no knowledge of the training corpus, but only knowledge of the configuration $(f_s,P,S)$ of the deployed forecaster, from which $\mathcal F_{\mathrm{lock}}$ follows by arithmetic. In the taxonomy of adversarial machine learning\autocite{vassilevAdversarialMachineLearning2024} this places the weakness with the availability violations, since what the model would lose at a candidate frequency would be lost for every input and every user, and it extends to the integrity violations of the evasion type once an adversary uses that loss to make a manipulation pass unnoticed;\footnote{The earlier deliverable classified the weakness as an availability violation only. That classification is kept for the weakness itself and extended here to its adversarial use, because the masking threat model developed in this section seeks to hide a manipulation from detection rather than to degrade the service, which the taxonomy associates with evasion and with the integrity of the output.} in both readings the attack lies outside the white-box assumption, since no access to the model is needed.

\autoref{tab:threatModel} reads the five elements of the adversarial-masking threat model\autocite{vassilevAdversarialMachineLearning2024}, separating the actor's capability and access path from the manipulation, consequence and exposure left after the proposed mitigation. Every row is conditional: it describes the exposure that would follow if the loss at the candidate sites were confirmed, and without that confirmation only the knowledge of where to look remains.

\begin{table}[!t]
    \centering
\small
    \setlength{\tabcolsep}{5pt}
    \begin{tabular}{@{}M{0.215\linewidth} M{0.735\linewidth}@{}}
        \toprule
        \textbf{Element} & \textbf{Reading for this system} \\
        \midrule
        Adversary capabilities & No model access of any kind and no knowledge of the training corpus, only the ability to drive a periodic load on the photovoltaic plant, which an inverter switching schedule or a switchable load already provides. \\
        \addlinespace[3pt]
        Attack surface & The raw telemetry stream before tokenisation, reached by acting on the physical process or on the data path. \\
        \addlinespace[3pt]
        Manipulation of telemetry & A small periodic component at a candidate frequency, which is legitimate plant behaviour and survives any plausibility check on the raw signal. \\
        \addlinespace[3pt]
        Operational consequence & The component is present in the plant and, if structural aliasing is confirmed, under-represented in the model, so dispatch and curtailment are decided on an incomplete picture, and a detector reading the forecast residual is looking through the same tokeniser that lost it. \\
        \addlinespace[3pt]
        Residual risk & The pre-registered mitigation raises overlap by shortening the stride at fixed patch size, so it thins the stride branch alone and $\mathcal{F}_{P}$ survives at every deployed configuration. Shrinking the patch instead thins $\mathcal{F}_{P}$, and with it $\mathcal{F}_{S}$ once $S \le P$ forces the stride down as well; but that raises both fundamentals rather than closing either comb, so the candidate set is relocated and thinned, not removed. \\
        \bottomrule
    \end{tabular}
\caption{The adversarial-masking threat model, read through the five elements of the NIST taxonomy of attacks on machine-learning systems\autocite{vassilevAdversarialMachineLearning2024}. The reading separates what the actor needs from what the actor changes, and both from the exposure that survives the mitigation this report pre-registers. Its distinguishing feature is the first row: the capability required is not access to the model but the ability to make the plant do something the plant is entitled to do, so nothing in the manipulation is anomalous at the point where a plausibility check would sit. \autoref{tab:riskRegister} carries this as its adversarial-masking row at high residual exposure.}
    \label{tab:threatModel}
\end{table}


This is the adversarial masking row of \autoref{tab:riskRegister}, carried at high residual: no control available to an operator closes it, since a spectral monitor placed after tokenisation would inherit the very blind spot the attacker exploits.

\subsection{Threat models beyond the adversarial one}

The framework asks for systematic vulnerability, not only for attacks, and the more likely exposure has no adversary in it. On one-minute telemetry, the arithmetic of \autoref{eq:branches} puts the patch branch of a $P{=}16$ tokeniser at periods of $16/k$ minutes, that is at $16$, $8$, $5.3$, $4$, $3.2$, $2.7$ and $2.3$ minutes. Those are the timescales of inverter duty cycling, of maximum-power-point dither and of cloud-shadow transients: ordinary plant behaviour, not an attack, sitting on the candidate frequencies at which the model may be least able to represent it. If the loss is confirmed, a monitoring system built on such a forecaster would be systematically less sensitive there, permanently and for reasons no operational log would reveal. \autoref{tab:riskRegister} records this as the systematic blind spot with no adversary, and it is why the candidate set belongs in the deployment documentation of such a system.

The third exposure, \autoref{tab:riskRegister}'s configuration drift, is governance rather than signal processing. Because the blind set is a function of configuration, a change that no one would route through a security review, a resampling of telemetry from one minute to thirty seconds, a swap of checkpoint, a patch size chosen for latency, relocates it. $H3a$ is the claim that would make this tractable: the relocation is predicted in closed form by $f_s/S$, and if the sites follow that prediction, recomputing the candidate set after a configuration change is arithmetic rather than re-measurement.

\subsection{Controls and residual exposure}

Two controls in this design are real, and the first answers \autoref{tab:riskRegister}'s defence placed inside the blind spot: detection has to sit before tokenisation, on the raw telemetry, because anything placed after it inherits the blind spot it is meant to guard; a spectral monitor on the raw stream, evaluated on the enumerated candidate set, is inexpensive precisely because that set is closed-form. And the semantic layer of \autoref{sec:appOntology} degrades actuation rather than failing silently: an assessment that is merely \valueof{PartiallyObservable} reduces the battery cap and lowers decision confidence, so uncertainty about the representation reaches the actuator instead of being resolved by assumption. That is the resilience the framework asks for, and it is the part of the design that does not depend on how the behavioural claim resolves.

One control cannot be credited. The pre-registered mitigation of \autoref{sec:agent}, increased patch overlap, is not established: its rule has two legs, both fixed in advance, and neither has been evaluated. The arithmetic bounds it in any case, since overlap moves the stride branch and leaves the patch branch untouched, so even a mitigation established on this design would be partial by construction. An operator who raises the overlap ratio and records the risk as mitigated has bought the false assurance of \autoref{tab:riskRegister}, the reason it is carried there at high residual exposure.
\endgroup

\subsection{Limitations}\label{sec:limitations}
The evidence has limits that no fit removes, and they are collected here rather than left to the appendices. Each geometry is one trained model from one seed, so a difference between two configurations is also a difference between two training runs, and the retrained models use half the published step budget (\autoref{sec:appRetraining}), which makes them comparable with one another but not with the released checkpoint. On the six geometries with $S\nmid P$ up to sixteen of the most recent context samples enter no token (\autoref{tab:sweepContext}), and these are the geometries that identify the stride branch. The token count moves with the stride, so the mitigation slope cannot separate overlap from sequence length. The $64$-sample horizon has a native resolution of $8$\,Hz, so a recovered line is located near a tone rather than resolved from its neighbours. The p32-s20 checkpoint behaves anomalously on the photovoltaic series and its training record is still to be checked. Finally, the measured series carries no energy on $\mathcal F_{\mathrm{lock}}$, so the domain informs the risk reading but tests none of the hypotheses. \textcolor{red}{\textbf{To be completed with the results.} A comparison with the findings of Pagani et al., placing the verdicts of \autoref{sec:bayesResults} beside theirs.}

